\subsection{Spectral high temperature: partial positive results and the remaining gap}
\label{ian:spectral-high-temperature}

Status: partial. Hereditary full AA for the entire spectral high-temperature class, including quenched SK at fixed sufficiently small positive temperature, is not established here. We prove hereditary simulation for every one-query acceptance function throughout that class, and hereditary full AA under an additional dimension-independent Frobenius bound. We also exclude the proposed Hadamard forrelation statistic, including its bounded polynomial amplifications, as an obstruction. The all-query positive theorem uses the uniform-AA import; the other results below do not.

Write
\[
 K=\beta J,\qquad
 \mu_{K,h}(x)=Z^{-1}\exp\left(\frac12x^\top Kx+h^\top x\right),
 \qquad K=K^\top,\quad K_{ii}=0.
\]
Throughout, full AA has the acceptance-simulation meaning of
Section~\ref{tht:target}. Computation and preprocessing are unrestricted.

\begin{proposition}[Pinning closure and precise non-coverage of SK]
\label{ian:spectral-pinning-and-noncoverage}
Pinning coordinates outside \(S\) to \(z\) produces the Ising law with
interaction matrix \(K_{SS}\) and field
\[
 h'_S=h_S+K_{S,S^c}z.
\]
Consequently, both conditions
\[
 \|K\|_{\mathrm{op}}\le c,\qquad \|K\|_{\mathrm F}\le B
\]
are preserved by pinning, with unchanged bounds.

For the Gaussian SK normalization
\(J_{ij}=g_{ij}/\sqrt n\), \(i<j\), the hypothesis of
Theorem~\ref{thm:resolution-dynamics-arbitrary-degree} fails eventually
almost surely at every fixed \(\beta>0\). Its maximal admissible inverse
temperature is
\[
 \frac{1}{8\max_i\sum_j|J_{ij}|}
 =\frac{1+o(1)}{8\sqrt{2/\pi}\sqrt n}
 \qquad\text{almost surely}.
\]
In contrast, since \(\|J\|_{\mathrm{op}}\to2\) almost surely (F\"uredi--Koml\'os; Bai--Yin; from literature, unverified), every fixed \(0<\beta<c/2\) satisfies \(\beta\|J\|_{\mathrm{op}}\le c\) eventually almost surely.
\end{proposition}

\begin{proof}
Expanding the Hamiltonian after pinning gives the displayed field.
A principal submatrix is an orthogonal compression, so its operator
norm cannot increase; its squared Frobenius norm is a sub-sum of that
of the original matrix.

For SK, put
\[
 R_i=\sum_j|J_{ij}|,\qquad \gamma=\sqrt{2/\pi}.
\]
Then \(\E R_i=\gamma(n-1)/\sqrt n\). As a function of its independent
Gaussian row entries, \(R_i\) is \(1\)-Lipschitz. Gaussian concentration
and a union bound give
\[
 \Pr\left[
 \max_i\left|R_i-\frac{\gamma(n-1)}{\sqrt n}\right|
 >\sqrt{6\log n}\right]\le 2n^{-2}.
\]
Borel--Cantelli yields
\[
 \max_iR_i=\gamma\sqrt n+O(\sqrt{\log n})
 \qquad\text{almost surely}.
\]
The pair-interaction hypothesis in
\texttt{high\_temperature/classical/05\_arbitrary\_degree.tex},
Theorem~\ref{thm:resolution-dynamics-arbitrary-degree}, is exactly
\(\beta\max_iR_i\le1/8\), proving the assertion.
Its convention \(\beta h_i\) instead of \(h_i\) is immaterial for
\(\beta>0\), since fields are arbitrary; at \(\beta=0\) the law is a
product law.
\end{proof}

Thus the existing weighted-interaction theorem does not cover fixed
positive-temperature SK. This non-coverage is a limitation of that
sufficient condition, not an AA lower bound.

\begin{lemma}[A hereditary spectral variance inequality]
\label{ian:spectral-poincare}
Suppose \(\|K\|_{\mathrm{op}}\le c<1/2\), and set
\(C_c=(1-2c)^{-1}\). For every real function \(f\),
\[
 \Var_{\mu_{K,h}}f
 \le C_c\sum_i
       \E_{\mu_{K,h}}\Var(f\mid X_{-i})
 \le C_c\sum_i I_i^{\mu_{K,h}}(f).
 \tag{ian-SP}
\]
The same constant works after every pinning and for every finite field.
In particular,
\[
 \Cov_{\mu_{K,h}}(X)\preceq C_c I.
\]
\end{lemma}

\begin{proof}
Choose \(\delta>0\) with \(2c+\delta<1\), and put
\[
 A=K+(c+\delta)I,\qquad \kappa=\|A\|_{\mathrm{op}}\le2c+\delta<1.
\]
Adding this diagonal changes the Hamiltonian only by a constant.
The Hubbard--Stratonovich joint law has latent marginal
\[
 \nu(dy)\propto
 \exp\left(-\frac12y^\top A^{-1}y\right)
 \prod_i2\cosh(h_i+y_i)\,dy,
\]
and \(X\mid Y=y\) has independent coordinates with means
\(\tanh(h_i+y_i)\). Denote this product law by \(P_y\).
The negative log density \(V\) of \(\nu\) satisfies
\[
 \nabla^2V(y)
 =A^{-1}-\diag\bigl(\sech^2(h_i+y_i)\bigr)
 \succeq(\kappa^{-1}-1)I.
\]
Hence
\[
 \Var_\nu g\le\frac{\kappa}{1-\kappa}\E_\nu\|\nabla g\|_2^2.
 \tag{ian-LP}
\]
For completeness, the strong-convexity variance bound follows from the
Langevin semigroup \(P_t\) with generator
\(\Delta-\nabla V\cdot\nabla\). Its flow Jacobian contracts by
\(e^{-\rho t}\) when \(\nabla^2V\succeq\rho I\), giving
\[
 \E_\nu\|\nabla P_tg\|^2
 \le e^{-2\rho t}\E_\nu\|\nabla g\|^2.
\]
Integration by parts and convergence to equilibrium yield
\[
 \Var_\nu g
 =2\int_0^\infty\E_\nu\|\nabla P_tg\|^2\,dt
 \le\rho^{-1}\E_\nu\|\nabla g\|^2.
\]
Here \(V\) has bounded Hessian and is strongly convex, so these identities
apply directly to smooth bounded Lipschitz \(g\); truncation gives the
usual extension.

Apply (ian-LP) to \(g(y)=\E_{P_y}f\). Differentiating gives
\[
 \partial_i g(y)=\Cov_{P_y}(f,X_i).
\]
The normalized centered coordinates under \(P_y\) are orthonormal.
Bessel's inequality, together with \(\Var_{P_y}X_i\le1\), gives
\[
 \sum_i\Cov_{P_y}(f,X_i)^2
 \le \Var_{P_y}f.
\]
Writing \(R=\E_\nu\Var_{P_y}f\), total variance therefore gives
\[
 \Var_\mu f=R+\Var_\nu g
 \le \frac{R}{1-\kappa}.
\]
Product variance tensorization and conditional total variance imply
\[
 R
 \le\sum_i\E\Var(f(X)\mid Y,X_{-i})
 \le\sum_i\E_\mu\Var(f(X)\mid X_{-i}).
\]
Letting \(\delta\downarrow0\) proves the first inequality in (ian-SP).
Since
\[
 \Var(f\mid X_{-i})
 =\Var(X_i\mid X_{-i})(D_if)^2\le(D_if)^2,
\]
the second follows. Apply it to linear functions to obtain the covariance
bound. Pinning closure was proved above.
\end{proof}

The covariance conclusion excludes a particular kind of two-component
structure, including after pinning. If an event has probability \(p\),
and the two conditional mean vectors are \(m_1,m_0\), then total
covariance gives
\[
 p(1-p)\|m_1-m_0\|_2^2\le C_c.
\]
Thus a balanced split cannot have conditional mean vectors separated
by order \(\sqrt n\). This does not exclude labels whose two conditional
laws have matching first moments.

\begin{theorem}[Hereditary AA for one quantum query]
\label{ian:one-query-spectral}
Under \(\|K\|_{\mathrm{op}}\le c<1/2\), every one-query quantum acceptance
function has a deterministic simulator of mean absolute error at most
\(\epsilon\), using at most
\[
 \min\left\{n,\left\lceil\frac{16C_c}{\epsilon^2}\right\rceil\right\}
\]
spin queries. The bound is uniform over \(K,h,n\), all pinnings, and
queryable auxiliary spins satisfying the same bound on the full
interaction matrix.
\end{theorem}

\begin{proof}
After purifying the algorithm and diagonalizing the response register
of its single oracle call, its state immediately after that call is
\[
 \psi_x=\psi_0+\sum_i x_i\psi_i.
\]
The vectors \(\psi_0,\psi_1,\ldots,\psi_n\) are mutually orthogonal and
their squared norms sum to one. The vector \(\psi_0\) includes the
oracle-invariant response subspace. There is an effect \(0\preceq E\preceq I\)
such that \(Q(x)=\langle\psi_x,E\psi_x\rangle\).
Put
\[
 p_i=\|\psi_i\|^2,\qquad a_i=\langle\psi_i,E\psi_i\rangle.
\]
These quantities are known before any input query, and
\(0\le a_i\le p_i\), \(\sum_i p_i\le1\).

Let \(S\) contain the \(k\) largest \(p_i\)'s. Unless all coordinates
have been included,
\[
 p_*:=\max_{i\notin S}p_i\le\frac1{k+1}.
\]
For every input,
\[
 D_iQ(x)=2\Re\langle\psi_i,E\psi_x\rangle-2x_i a_i.
\]
Orthogonality and \(\|E\|\le1\) give
\[
 \begin{split}
 \sum_{i\notin S}(D_iQ(x))^2
 &\le8\sum_{i\notin S}|\langle\psi_i,E\psi_x\rangle|^2
       +8\sum_{i\notin S}a_i^2\\
 &\le8p_*\|E\psi_x\|^2+8p_*\sum_i p_i
 \le16p_*.
 \end{split}
\]
After querying \(X_S\), apply Lemma~\ref{ian:spectral-poincare} on every
conditional law:
\[
 \Var(Q\mid X_S)\le\frac{16C_c}{k+1}.
\]
Output \(D=\E[Q\mid X_S]\). Then
\[
 \E|Q-D|
 \le\sqrt{\E\Var(Q\mid X_S)}
 \le4\sqrt{\frac{C_c}{k+1}}.
\]
The stated choice of \(k\), or reading the entire input, proves the
claim. The same argument begins on any pinned law.
\end{proof}

This proof replaces posterior access by direct revelation of the
algorithm's largest input-independent query weights. It excludes every
one-query full-AA obstruction in the spectral class, not merely the
specific Hadamard witness.

\begin{lemma}[Centered quadratic exponential moments]
\label{ian:centered-quadratic-moment}
Let \(\xi_i\) be independent centered random variables satisfying
\(\E e^{s\xi_i}\le e^{s^2/2}\) for all real \(s\).
For a symmetric zero-diagonal matrix \(K\), put
\(F=\frac12\xi^\top K\xi\). Whenever
\(2|t|\|K\|_{\mathrm{op}}<1\),
\[
 \E e^{tF}
 \le
 \exp\left(
 \frac{t^2\|K\|_{\mathrm F}^2}
      {1-4t^2\|K\|_{\mathrm{op}}^2}
 \right).
 \tag{ian-QM}
\]
\end{lemma}

\begin{proof}
Independently choose a random subset \(S\) by fair coin flips, and let
\[
 C_S=\xi_S^\top K_{S,S^c}\xi_{S^c}.
\]
Conditional on \(\xi\), \(\E_S C_S=F/2\), so Jensen gives
\[
 \E e^{tF}\le\E_S\E_\xi e^{2tC_S}.
\]
Fix \(S\) and write \(M=K_{S,S^c}\). Conditioning on \(\xi_S\), applying
the assumed linear moment bound, and then introducing a standard
Gaussian vector \(G\), gives
\[
 \begin{split}
 \E e^{2t\xi_S^\top M\xi_{S^c}}
 &\le\E_{\xi_S}e^{2t^2\|M^\top\xi_S\|^2}\\
 &=\E_{\xi_S,G}e^{2t\xi_S^\top MG}\\
 &\le\E_G e^{2t^2\|MG\|^2}
 =\det(I-4t^2M^\top M)^{-1/2}.
 \end{split}
\]
Using \(-\log(1-u)\le u/(1-u_{\max})\),
\[
 \log\E e^{2tC_S}
 \le\frac{2t^2\|M\|_{\mathrm F}^2}
          {1-4t^2\|K\|_{\mathrm{op}}^2}
 \le\frac{t^2\|K\|_{\mathrm F}^2}
          {1-4t^2\|K\|_{\mathrm{op}}^2}.
\]
Here \(\|M\|_{\mathrm{op}}\le\|K\|_{\mathrm{op}}\), and symmetry gives
\(2\|M\|_{\mathrm F}^2\le\|K\|_{\mathrm F}^2\).
The bound is independent of \(S\).
\end{proof}

\begin{theorem}[Hereditary full AA with a fixed Frobenius budget]
\label{ian:frobenius-full-aa}
Assume the uniform-AA import, in the simulation form \(s_U(T,e)=\operatorname{poly}(T,1/e)\) used in \texttt{foundations/transfers/01\_posterior\_access.tex} (Proposition~\ref{prop:threshold-local-posterior}) and derived from the inequality of Section~\ref{ht:conventions-and-imported-uniform-theorem}, through the product simulation theorem \texttt{simulation/positive\_examples/02\_arbitrary\_blocks.tex}, Theorem~\ref{thm:threshold-block-products}.
For every fixed finite \(B\), the class
\[
 \|K\|_{\mathrm{op}}\le\frac18,\qquad
 \|K\|_{\mathrm F}\le B,\qquad h\in\R^n
 \tag{ian-FB}
\]
has hereditary full AA, uniformly over \(n,K,h\).

More precisely, let
\(s_{\mathrm{prod}}(T,\eta)=S_B(T,\eta;1)\) be the arbitrary-product
query bound from that theorem. A valid bound here is
\[
 s_{\mathrm{prod}}\left(
 T,\epsilon^2e^{-16B^2/3}\right).
 \tag{ian-FQ}
\]
For fixed \(B\), this is a polynomial in \(T,1/\epsilon\), independent
of dimension and fields.
\end{theorem}

\begin{proof}
The map \(m\mapsto\tanh(h+Km)\) is a strict Euclidean contraction of
\([-1,1]^n\). Let \(m\in(-1,1)^n\) be its unique fixed point, and let
\(P_m\) be the product-spin law with means \(m_i\). Thus
\[
 \operatorname{arctanh}(m)=h+Km.
\]
Writing \(\xi=X-m\), direct comparison of the two densities gives
\[
 \frac{d\mu_{K,h}}{dP_m}(x)
 =\frac{e^{F(x)}}{\E_{P_m}e^F},
 \qquad F(x)=\frac12(x-m)^\top K(x-m).
 \tag{ian-centered-density}
\]
The centered product spins satisfy
\(\E e^{s\xi_i}\le e^{s^2/2}\): the second derivative of their log
moment-generating function is a tilted spin variance, at most one.
Also \(\E_{P_m}F=0\), because \(K\) has zero diagonal.
Consequently \(\E_{P_m}e^F\ge1\).

Apply Lemma~\ref{ian:centered-quadratic-moment} at \(t=2\).
Under (ian-FB),
\[
 \E_{P_m}e^{2F}
 \le\exp\left(\frac{4B^2}{1-16/64}\right)
 =e^{16B^2/3}.
\]
For \(r=d\mu/dP_m\), this proves
\[
 \|r\|_{L^2(P_m)}\le e^{8B^2/3}=:M_B.
\]
Use product AA to obtain \(D\in[0,1]\) with
\[
 \E_{P_m}|D-Q|\le\eta,
 \qquad \eta=\epsilon^2M_B^{-2}.
\]
Cauchy--Schwarz and \(|D-Q|^2\le|D-Q|\) give
\[
 \E_\mu|D-Q|
 \le M_B\bigl(\E_{P_m}|D-Q|^2\bigr)^{1/2}
 \le M_B\sqrt\eta=\epsilon.
\]
The simulator directly queries the physical input, so its query count
is unchanged. Pinning preserves both matrix bounds; its arbitrarily
large induced fields are handled by recomputing the fixed point.
\end{proof}

Here bounded density moments relative to a suitable product law replace
local posterior access entirely. No latent field is sampled.

In particular, fixed-rank spectral families satisfy this theorem:
if \(\rank K\le r\) and \(\|K\|_{\mathrm{op}}\le1/8\), then
\(\|K\|_{\mathrm F}\le\sqrt r/8\). The Frobenius restriction is only
sufficient. For example, arbitrarily many independent weak edges violate
a fixed Frobenius budget while retaining full AA by
Theorem~\ref{thm:threshold-block-products}.

This theorem does not include fixed-temperature SK:
\[
 \|J\|_{\mathrm F}^2=\frac2n\sum_{i<j}g_{ij}^2,
 \qquad
 \frac{\|J\|_{\mathrm F}^2}{n}\longrightarrow1
 \quad\text{in probability}.
\]
The last assertion follows immediately from its mean and variance.
Thus the density-transfer cost above is exponential in order
\(\beta^2n\), rather than dimension independent.

\begin{proposition}[Hadamard forrelation concentrates under every pinning]
\label{ian:hadamard-concentration}
Let \(O\) be any real orthogonal \(N\times N\) matrix, and define
\[
 R=\begin{pmatrix}0&O\\O^\top&0\end{pmatrix},
 \qquad K=\lambda R,\qquad |\lambda|\le c<1/2.
\]
For \(x=(u,v)\in\{-1,1\}^{2N}\), put
\[
 \phi(u,v)=\frac{u^\top Ov}{N},
 \qquad Q(u,v)=\frac{1+\phi(u,v)}2.
\]
Then \(Q\) is a one-query acceptance function and, for every field and
every pinning,
\[
 \Var(\phi)\le\frac{2C_c}{N},
 \qquad
 \Var(Q)\le\frac{C_c}{2N}.
 \tag{ian-FC}
\]
These inequalities also hold with additional queryable auxiliary spins
if the full interaction matrix has operator norm at most \(c\).

Consequently the zero-query output \(D=\E Q\), computed for the specified
law and pinning, has error at most \(\sqrt{C_c/(2N)}\).
\end{proposition}

\begin{proof}
Orthogonality gives \(|\phi|\le1\), \(R=R^\top\), and \(R^2=I\).
A phase query prepares
\[
 \psi_x=\frac1{\sqrt{2N}}\sum_{i=1}^{2N}x_i|i\rangle.
\]
Measuring the projection \((I+R)/2\) accepts with probability
\[
 \left\langle\psi_x,\frac{I+R}{2}\psi_x\right\rangle
 =\frac12+\frac{x^\top Rx}{4N}=Q(x).
\]
Moreover,
\[
 D_{u_i}\phi=\frac{(Ov)_i}{N},\qquad
 D_{v_j}\phi=\frac{(O^\top u)_j}{N},
\]
and therefore, pointwise on the entire cube,
\[
 \sum_i(D_i\phi)^2
 =\frac{\|Ov\|^2+\|O^\top u\|^2}{N^2}
 =\frac2N.
\]
After pinning, the sum over unpinned coordinates can only decrease.
Auxiliary-spin derivatives vanish. Apply
Lemma~\ref{ian:spectral-poincare}; the mean absolute error assertion
follows from Cauchy--Schwarz.
\end{proof}

In particular,
\[
 \Pr\bigl[|\phi-\E\phi|\ge t\bigr]\le\frac{2C_c}{Nt^2}
\]
uniformly under pinnings. Such pinnings cannot create two components
with appreciable mass and a constant separation in this statistic.

\begin{corollary}[Polynomial amplifications of the statistic are simulable]
\label{ian:hadamard-polynomial-simulation}
In the preceding setting, let \(p:[-1,1]\to[0,1]\) be a polynomial of
degree \(d\ge1\). The function \(p(\phi)\) has a deterministic simulator
of mean error at most \(\epsilon\), using
\[
 O(C_c d^4\epsilon^{-2})
\]
visible-spin queries, uniformly in fields, pinnings, and dimension.
No quantum-realizability assumption on \(p(\phi)\) is needed.
\end{corollary}

\begin{proof}
Markov's polynomial inequality gives \(\|p'\|_\infty\le d^2\).
Using two independent copies to express variance,
\[
 \Var(p(\phi))\le d^4\Var(\phi)\le\frac{2C_cd^4}{N}.
\]
If \(N\ge2C_cd^4/\epsilon^2\), output \(\E p(\phi)\) without querying.
Otherwise read the at most \(2N\) unpinned visible spins and evaluate the
function exactly.
\end{proof}

For the normalized Walsh Hadamard matrix, \(O_{ij}=\pm N^{-1/2}\),
this is exactly the proposed spectral geometry:
\[
 \|R\|_{\mathrm{op}}=1,\qquad
 \max_i\sum_j|R_{ij}|=\sqrt N,\qquad
 \|\lambda R\|_{\mathrm F}^2=2\lambda^2N.
\]
Thus the forrelation calculation applies precisely where the existing
absolute-row-sum theorem ceases to apply. It rules out this statistic
and its polynomial amplifications, rather than proving full AA for
every acceptance function under the Hadamard law.

Replacing the law by a mixture of the \(+\lambda R\) and
\(-\lambda R\) laws changes the problem. The formal sign-spin
implementation
\[
 \exp\bigl(\lambda s\,u^\top Ov\bigr)
\]
has cubic interactions \(s u_i v_j\), so it is not the requested
pairwise Hamiltonian. Also, when \(s\) is queryable, querying it and
outputting \(\E[Q\mid s]\) simulates this particular witness to
\(O_c(N^{-1/2})\) error, by (ian-FC) applied within each sign component.
No pairwise auxiliary-spin realization preserving a separation is
constructed here.

\textbf{The remaining step.}
The spectral variance inequality does not supply polynomial reveal
progress for arbitrary multi-query acceptance functions. The one-query
proof uses orthogonal vectors with input-independent weights \(p_i\);
later quantum query states depend on the unknown input, and that
argument does not extend as written.

Precisely, what remains is a derivation, from a small spectral norm
alone, of the quantities \(K(T,1/a)\) and \(g(T,a)\) in
Proposition~\ref{tht:reveal-prop}, uniformly over all fields and pinnings,
with polynomial bounds. Alternatively, one needs a different
representation satisfying
Proposition~\ref{prop:threshold-local-posterior}, or a genuine
multi-query counterexample. The next subsection proves that canonical
Gaussian-coordinate posterior access and uniform local heat-bath
access cannot provide the required dimension-independent primitive.
It does not rule out other representations.
